\chapter{Geltungsbereich und Notation}
\label{b21guide:ch:foundations}

Dieser Leitfaden ergänzt Band~21 durch eine ausgewählte Übersicht
wichtiger Definitionen, Charakterisierungen und Folgerungen über
\emph{Folgen}. Er enthält keine Beweise oder Beweisskizzen.
Die Aussage-Nummern mit Präfix $21H$ gehören nur zu dieser Ergänzung.
Das Format folgt den Leitfäden zu den Bänden~3, 4, 5, 28 und~37.

\section{Voraussetzungen und Abgrenzung}

Aufgenommen werden Aussagen, in denen die endliche oder natürliche
Indexfolge, das Verhalten ab einem hinreichend großen Index oder
der Grenzübergang wesentlich ist. Eine bloße Umschreibung eines
allgemeinen Funktionssatzes in Folgennotation ist kein eigener
Gegenstand dieses Leitfadens.

Wie im Hauptband gibt es zwei Ebenen: Folgen in einer beliebigen
Menge $M$ und \emph{reelle Folgen}. Für die erste Ebene wird keine
Struktur auf $M$ vorausgesetzt. Ab Kapitel~\ref{b21guide:ch:limits}
ist die Zielmenge ausdrücklich $\mathbb R$ mit der in Band~19
entwickelten Arithmetik, Ordnung und Vollständigkeit.
Grenzwerte und Cauchy-Eigenschaft folgen nicht schon aus dem
bloßen Begriff einer Abbildung $\mathbb N\to M$.
Die benötigten Eigenschaften der reellen Zahlen werden hier
verwendet, ihre Theorie wird nicht erneut aufgebaut.

Allgemeine Funktionen und besondere Funktionseigenschaften
(Injektivität, Surjektivität, Bijektivität) bleiben in den
Bänden~5--8. Äquivalenzrelationen, Halbordnungen, totale Ordnungen
und allgemeine Resultate über Infima und Suprema bleiben in den
Bänden~11--15. Es werden weder beliebige geordnete Zielmengen noch
algebraische, metrische oder topologische Strukturen eingeführt.
Insbesondere bleibt die Verallgemeinerung auf metrische Räume
in Band~47. Reihen, Funktionenfolgen, Netze, Filter und
die algebraische Theorie endlicher Wörter sind ebenfalls nicht
Gegenstand dieses Leitfadens.

Beschränktheit, Monotonie und das Nichtverschwinden eines Nenners
sind, wo sie auftreten, ausdrücklich genannte Voraussetzungen
des einzelnen Satzes. Sie werden nicht für alle Folgen angenommen.
Als mengentheoretischer Hintergrund genügt ZF; ein zusätzliches
Auswahlaxiom wird für die aufgeführten Aussagen nicht benötigt.
Die Zusammenstellung ist eine Auswahl zentraler bekannter Resultate;
sie beansprucht keine vollständige Aufzählung aller bekannten
oder ableitbaren Sätze und keine historischen Erstzuordnungen.

\section{Folge und Indexkonvention}

\begin{definition}[Indizierte Familie und Folge]
\label{b21guide:def:sequence}
Für eine Funktion $a\colon I\to M$ bedeuten
$(a_i)_{i\in I}=a$ und $a_i=a(i)$ die Familiennotation.
Eine \emph{unendliche Folge in $M$} hat Indexmenge
$\mathbb N=\{0,1,2,\ldots\}$.
Für $r\in\mathbb N$ steht
$\mathbb N_{\geq r}=\{n\in\mathbb N:n\geq r\}$.
Eine bei $r$ beginnende Folge hat diese Indexmenge.
\end{definition}

\begin{definition}[Endliche Folge und Länge]
\label{b21guide:def:finite}
Für $m\in\mathbb N$ sei
\[
 [m]_0=\{0,\ldots,m-1\},\qquad [0]_0=\varnothing.
\]
Eine \emph{endliche Folge der Länge $m$} in $M$ ist eine Funktion
$a\colon[m]_0\to M$, geschrieben $(a_0,\ldots,a_{m-1})$.
Die Länge wird mit $\ell(a)$ bezeichnet.
Die Folge der Länge $0$ heißt \emph{leere Folge}.
\end{definition}

Eine Folge enthält die Zuordnung zu den Indizes; die Wertemenge
$\{a_n:n\in\mathbb N\}$ allein beschreibt sie nicht.
Gleiche Werte an verschiedenen Stellen sind erlaubt.
Die Zielmenge $M$ muß nicht vollständig durchlaufen werden.
Diese Hinweise sind Notationskonventionen, keine Wiederholung
allgemeiner Sätze über Funktionen.

\section{Wegweiser}

Kapitel~\ref{b21guide:ch:finite} enthält endliche Folgen,
Anfangsstücke und Rekursion.
Kapitel~\ref{b21guide:ch:subsequences} ordnet Restfolgen,
Teilfolgen, Aussagen über hinreichend große Indizes und Periodenstruktur.
Die Kapitel~\ref{b21guide:ch:limits} und~\ref{b21guide:ch:calculus}
sammeln Konvergenzkriterien und Grenzwertregeln für reelle Folgen;
Kapitel~\ref{b21guide:ch:calculus} enthält außerdem die endliche
Newtondarstellung, endliche Summationsformeln und die
Charakterisierung polynomialer Folgen.
Kapitel~\ref{b21guide:ch:compactness} behandelt Häufungswerte,
monotone Teilfolgen, Teilfolgenkriterien, Cauchy-Folgen,
quantitative Fehlerabschätzungen
und stabile Rekursionen.
Kapitel~\ref{b21guide:ch:monotone} enthält monotone Konvergenz
und bestimmte Divergenz.
Kapitel~\ref{b21guide:ch:standard} bietet gebräuchliche Grenzwerte,
Wurzelkriterien, Feketes Grenzwertsatz, Mittelwerte, Umkehrkriterien
und beide Fassungen
des Stolz--Cesàro-Kriteriums.
